\chapter{A1: Per-Step Rao--Blackwellisation --- Full Investigation}
\label{app:a1}

This appendix backs the A1 autopsy of \Cref{sec:autopsies} with its full evidence. A1 is the most
instructive of the negative results: the hypothesis it began with was wrong, the corrected diagnosis
identified a \emph{real} variance lever, and that lever was still declined---because it pays off only
in the regime the renderer already dominates. The episode is recorded here in detail both as the
substantiation of the summary and as a worked example of the measure-before-rewriting discipline that
governs \Cref{ch:optimization}.

\section{Premise}
\label{sec:a1-premise}

Profiling against the Mitsuba reference found the renderer roughly $2.85\times$ noisier per sample on
the scattering cloud under flat lighting. The first hypothesis named a concrete cause: Mitsuba folds
the analytic segment transmittance into the path throughput at \emph{every} bounce
($\beta \mathrel{*}= T_{\text{seg}}$), whereas this renderer only does so at the primary vertex (its
analytic-direct path). The proposed fix---``per-step Rao--Blackwellisation''---was to generalise that
fold to all bounces. The hypothesis was explicitly flagged as \emph{unprofiled}, and so was tested
before any change to the validated core estimator.

\section{Diagnosis}
\label{sec:a1-diagnosis}

A depth sweep on the cleanest possible scene---a single Gaussian (no overlap), constant environment,
$\alpha = 0.9$---isolates the effect. \Cref{tab:a1-depth} reports per-seed noise (standard deviation
across seeds) against the maximum path depth.

\begin{table}[htp]
  \centering
  \begin{tabular}{@{}rrrr@{}}
    \toprule
    Max depth & This renderer & Mitsuba (analog) & Ratio \\
    \midrule
    1 & 0.00599 & 0.00375 & $1.60\times$ \\
    2 & 0.00639 & 0.00147 & $4.34\times$ \\
    4 & 0.00643 & 0.00128 & $5.02\times$ \\
    8 & 0.00643 & 0.00128 & $5.02\times$ \\
    \bottomrule
  \end{tabular}
  \caption{Per-seed noise versus maximum path depth, single Gaussian under a constant environment
  ($\alpha = 0.9$, six seeds). This renderer's noise is \emph{flat} with depth; Mitsuba's
  \emph{drops}. The means agree, so both estimators are unbiased---the difference is purely variance,
  and it is set by the multiple-scatter term, not the primary vertex.}
  \label{tab:a1-depth}
\end{table}

Two control experiments place the gap. Disabling multiple-importance sampling (next-event estimation
only) changes the noise by just $1.06\times$, so it is not an environment-sampling artefact; and the
single-Gaussian scene reproduces the full gap, so it is not an effect of the argmin overlap region.
The gap lives entirely in the converged multiple-scatter term.

The corrected root cause emerged from reading Mitsuba's segment sampler rather than guessing at it.
The two renderers use \emph{different classes} of estimator. Mitsuba folds the analytic segment
transmittance into throughput at every segment---a \emph{track-length} (expected-value) estimator, in
which every path contributes a smoothly $T$-weighted escape. The architecture of this thesis, by
construction, makes a binary scatter-versus-escape decision per bounce (\Cref{sec:adt}): the
transmittance is \emph{consumed by that decision} and never folded into the throughput. This is a
\emph{collision} estimator. For transmittance-weighted estimands a track-length estimator has the
lower variance, and the gap widens with depth: in a near-conservative constant medium almost every
analog path eventually escapes to the same near-constant environment, so Mitsuba's walk self-averages
and its variance vanishes with depth, while the collision estimator's stays flat. This both
\emph{confirms} the original premise---it was the binary continuation, not the primary vertex, that
opened the gap, exactly as the depth-1 row (where the analytic-direct fold is already active) shows---
and \emph{corrects} an intermediate reading that had mistaken the two estimators for identical.

\section{Why the natural fix raises variance}
\label{sec:a1-fix}

A tempting alternative---adding the continuation ray's escape as a second sampling strategy, in the
textbook way a BSDF sample can hit an emitter---is \emph{strictly worse} than what the renderer
already does. The existing phase-importance-sampled next-event estimator evaluates the same direct
contribution but with an \emph{analytic} shadow-ray transmittance; the continuation strategy would
use a \emph{stochastic} escape indicator instead. Same expectation, higher variance. That rewrite
would raise the variance of the direct term, not lower it.

\section{Conclusion}
\label{sec:a1-conclusion}

The corrected diagnosis identifies a genuine lever: granting the renderer a track-length throughput
would recover the missing self-averaging under flat lighting. It was nonetheless declined, for two
reasons. First, the lever helps only the flat- or constant-environment regime; on the environment-map
showcase the renderer already matches and beats Mitsuba (\Cref{ch:results}), cleaner in total noise
and free of its fireflies, because multiple-importance sampling closes the variance there. Second,
folding transmittance into the throughput at every bounce means abandoning the binary scatter/escape
decision that \emph{is} the sort-free argmin scheme---a rewrite of the validated core estimator that
would undo the architecture's central contribution to win a regime that is not the target. The honest
verdict is therefore a characterised trade-off rather than a missed optimisation: no single estimator
beats both, and the one this renderer chose is the one its showcase rewards. A hybrid that keeps the
argmin traversal but folds the analytic transmittance into throughput---generalising the primary-vertex
fold to all bounces without discarding the single-trace structure---remains the open direction should
the flat-environment gap ever become worth closing.

A small mean residual, separate from the variance question above, persists at convergence: under the
single-Gaussian constant-environment test the two renderers' means differ by about $0.5\%$ (this
renderer slightly brighter). It is consistent with the low-density interior and coloured-albedo
residuals noted in \Cref{ch:validation}, and a likely cause is the differing minimum-throughput cull
threshold. It is tracked as an open item and is unrelated to A1.
